% !TeX root = Report.tex
\chapter*{Заключение}
\addcontentsline{toc}{chapter}{Заключение}
\enlargethispage{2\baselineskip}

В работе рассматривалась задача предобработки изображения перед применением
преобразования Хафа для обнаружения прямых линий. Основное внимание было
уделено точности оценивания параметров прямых $(\rho,\theta)$. Полученные
результаты показывают, что качество преобразования Хафа после предобработки
определяется не только степенью подавления шума, но и тем, какие активные
пиксели остаются после обработки изображения.

В ходе работы была построена модель изображения с одной или несколькими
дискретными прямыми и описана связь этой модели с комплексными рядами,
возникающими после DFT. Для варианта \texttt{CSSA ROW.ROW} показано, что в
случае одной прямой строки после DFT имеют низкоранговую экспоненциальную
структуру. Для варианта \texttt{CSSA COL.ROW} получены оценки ранга траекторной
матрицы, позволяющие выбирать число сохраняемых компонент с учетом структуры
прямой, а не только эмпирически.

Отдельно исследована проблема растекания после обратного DFT. Показано, что
она существенно влияет на преобразование Хафа: даже при малой ошибке
восстановления лишние активные пиксели могут создавать побочные максимумы в
пространстве параметров. Для уменьшения этого эффекта рассмотрены два подхода:
пороговая фильтрация по оценке остаточного шума и коррекция частот с помощью
ESPRIT.

Численные эксперименты показали, что для одной прямой методы
\texttt{CSSA ROW.ROW}, \texttt{CSSA COL.ROW}, \texttt{MAX.ROW} и Wiener в
ряде конфигураций достигают ошибки, близкой к нижнему уровню, заданному
дискретизацией преобразования Хафа. При переходе к двум прямым более устойчивыми
оказываются методы, контролирующие число активных пикселей: \texttt{MAX.ROW}
и \texttt{CSSA ROW.ROW + ESPRIT}. При этом \texttt{MAX.ROW} использует
априорное знание числа прямых, тогда как подход на основе CSSA дополнительно дает
остаточную матрицу, по которой можно оценивать уровень шума и выбирать порог.

Таким образом, в работе показано, что CSSA может быть эффективным инструментом
предобработки перед преобразованием Хафа, но его применение требует учета
ранговой структуры сигнала, выбора числа компонент и контроля растекания.
Дальнейшее развитие работы связано с автоматическим выбором параметров CSSA,
анализом устойчивости ESPRIT для близких прямых и расширением экспериментов на
другие шумовые модели. Все файлы выложены на GitHub:
\href{https://github.com/Programmer-Alexey/SSA-in-Image-processing}{SSA in
Image Processing}.
